\chapter[14]{Local Derivative Estimates for Curvature}
\label{ch:chow-ii-local-derivative-estimates}

Let \(g(t)\) be a Ricci flow.  Norms, covariant derivatives, and balls use
the indicated metric; \(\operatorname{Rm}\) and \(\operatorname{Rc}\) denote
its Riemann and Ricci tensors.

\section{Metric comparison and cutoff functions}

\begin{lemma}[Ricci curvature bounds change the metric]
\label{lem:chow-ii-ricci-bounds-metric-comparison}
\source{chow-et-al-ricci-flow-part-ii:page-0256}
\dcref{ch14:14.1}
\group{Metric Comparison}

Let \(g(t)\), \(t\in[0,T)\), be a Ricci flow and let \(\mathcal U\) be an
open subset of the manifold.
\begin{enumerate}
\item If \(\operatorname{Rc}\le K_1g\) on \(\mathcal U\times[0,T)\), then
\[
g(x,t)\ge e^{-2K_1T}g(x,0).
\]
\item If \(\operatorname{Rc}\ge-K_2g\) on \(\mathcal U\times[0,T)\), then
\[
g(x,t)\le e^{2K_2T}g(x,0).
\]
\end{enumerate}
\end{lemma}
\begin{proof}
\sketch For \(0\ne V\in T_xM\), the Ricci flow equation gives
\[
\log\frac{g(t)(V,V)}{g(0)(V,V)}
=-2\int_0^t\frac{\operatorname{Rc}_{g(s)}(V,V)}{g(s)(V,V)}\,ds.
\]
In the first case the integrand is at most \(K_1\), and in the second it is
at least \(-K_2\).  Integrating and exponentiating proves both inequalities.
\end{proof}

\begin{remark}[Uniform equivalence under two-sided Ricci bounds]
\label{rem:chow-ii-uniform-metric-equivalence}
\source{chow-et-al-ricci-flow-part-ii:page-0256}
\uses{lem:chow-ii-ricci-bounds-metric-comparison}
\dcref{ch14:14.2}
\group{Metric Comparison}


If \(-K_2g\le\operatorname{Rc}\le K_1g\), then
\[
e^{-2K_1T}g(0)\le g(t)\le e^{2K_2T}g(0)
\]
on the common space-time domain.
\end{remark}

\begin{lemma}[Gradient and Laplacian estimates for a cutoff function]
\label{lem:chow-ii-cutoff-gradient-laplacian-estimates}
\source{chow-et-al-ricci-flow-part-ii:page-0257,chow-et-al-ricci-flow-part-ii:page-0258,chow-et-al-ricci-flow-part-ii:page-0259}
\uses{lem:chow-ii-ricci-bounds-metric-comparison}
\dcref{ch14:14.3}
\group{Cutoff Functions}


Suppose \(\overline B_{g(0)}(p,r)\Subset\mathcal U\),
\(|\operatorname{Rm}|\le K\) on \(B_{g(0)}(p,r)\times[0,\tau]\), and
\(\tau\le\alpha/K\).  Put
\(\Theta(x,t)=tK^2|\nabla\operatorname{Rc}(x,t)|^2\).  There is a Lipschitz cutoff
\(\eta:\mathcal U\to[0,1]\), supported in \(B_{g(0)}(p,r)\) and equal to
one on \(B_{g(0)}(p,r/2)\), such that, in the barrier sense,
\[
|\nabla\eta|_{g(t)}^2\le\frac{C(\alpha,n)}{r^2}\eta
\]
and
\[
-\Delta_{g(t)}\eta\le\frac{C(\alpha,\sqrt K r,n)}{r^2}
+\frac{C(\alpha,n)}{K^{3/2}r}
 \sup_{0\le s\le t}(\eta\Theta)^{1/2}(x,s).
\]
\end{lemma}
\begin{proof}
\sketch Choose a nonincreasing smooth \(\lambda:[0,\infty)\to[0,1]\), equal to one
on \([0,1/2]\) and zero on \([7/8,\infty)\), and set
\(\kappa(s)=\lambda(s/r)^2\).  It may be chosen so that
\[
0\le-\kappa'\le 6r^{-1}\sqrt\kappa,
\qquad |\kappa''|\le250r^{-2}.
\]
Let \(d_0=d_{g(0)}(p,\cdot)\) and \(\eta=\kappa\circ d_0\).  Metric
comparison, using \(|\operatorname{Rc}|\le C(n)K\), gives
\(|\nabla\eta|_{g(t)}^2\le C(\alpha,n)r^{-2}\eta\).

Hessian comparison for \(g(0)\), whose sectional curvature is bounded below
by \(-K\), gives on the support of \(\kappa'\)
\[
\nabla^0\nabla^0\eta\ge-C(\sqrt K r)r^{-2}g(0).
\]
Furthermore
\[
\Delta_{g(t)}\eta
=g(t)^{ij}\nabla_i^0\nabla_j^0\eta
+g(t)^{ij}(\Gamma(0)-\Gamma(t))^k_{ij}\nabla_k\eta.
\]
Since \(\partial_t\Gamma=\nabla\operatorname{Rc}*1\), metric comparison and the gradient
bound imply
\[
\left|g(t)^{ij}(\Gamma(0)-\Gamma(t))^k_{ij}\nabla_k\eta\right|
\le\frac{C(\alpha,n)}{Kr}
\int_0^t\eta^{1/2}|\nabla\operatorname{Rc}|(x,s)\,ds.
\]
Writing \(|\nabla\operatorname{Rc}|=K^{-1}s^{-1/2}\Theta^{1/2}\) and using
\(t\le\alpha/K\) bounds this by the second term in the asserted Laplacian
estimate.  The first term follows from Hessian comparison and uniform metric
equivalence.  Calabi's support-function argument supplies the same bounds at
the cut locus.
\end{proof}

\begin{lemma}[Cutoff functions with bounded derivatives]
\label{lem:chow-ii-cutoff-bounded-derivatives}
\source{chow-et-al-ricci-flow-part-ii:page-0259}
\uses{lem:chow-ii-cutoff-gradient-laplacian-estimates}
\dcref{ch14:14.4}
\group{Cutoff Functions}


Under the hypotheses of Lemma~\ref{lem:chow-ii-cutoff-gradient-laplacian-estimates},
suppose additionally that
\[
tK^2|\nabla\operatorname{Rc}|^2\le C(\alpha,r,n)K^4
\]
on \(B_{g(0)}(p,r/2)\times[0,\tau]\).  For every integer \(m\ge2\) there is
a cutoff \(\eta\), supported in \(B_{g(0)}(p,r/2^m)\) and equal to one on
\(B_{g(0)}(p,r/2^{m+1})\), such that
\[
|\nabla\eta|^2\le C(\alpha,m,n)r^{-2}\eta,
\qquad -\Delta\eta\le C(\alpha,K,r,m,n).
\]
\end{lemma}
\begin{proof}
\sketch Apply Lemma~\ref{lem:chow-ii-cutoff-gradient-laplacian-estimates} with the
outer radius \(r/2^m\).  The assumed bound on \(\Theta\), together with
\(0\le\eta\le1\), turns its final supremum into a constant depending only on
the displayed parameters.  Absorbing the powers of \(2^m\) into the constant
gives both estimates.
\end{proof}

\section{Global derivative estimates}

\begin{theorem}[Global derivative estimates]
\label{thm:chow-ii-global-curvature-derivative-estimates}
\source{chow-et-al-ricci-flow-part-ii:page-0259,chow-et-al-ricci-flow-part-ii:page-0260,chow-et-al-ricci-flow-part-ii:page-0261}
\dcref{ch14:14.5}
\group{Curvature Derivative Estimates}

For \(\alpha>0\) and \(m,n\in\mathbb N\), there is a constant
\[
C=C(m,n,\max\{\alpha,1\})
\]
with the following property.  If the maximum principle applies to a Ricci
flow \(g(t)\) and
\[
|\operatorname{Rm}(x,t)|\le K\qquad(0\le t\le\alpha/K),
\]
then
\[
|\nabla^m\operatorname{Rm}(x,t)|\le \frac{CK}{t^{m/2}}
\qquad(0<t\le\alpha/K).
\]
\end{theorem}
\begin{proof}
\sketch The curvature evolution inequalities have the form
\[
(\partial_t-\Delta)|\nabla^j\operatorname{Rm}|^2
\le-2|\nabla^{j+1}\operatorname{Rm}|^2
+C(j,n)\sum_{i=0}^j|\nabla^i\operatorname{Rm}|,|\nabla^{j-i}\operatorname{Rm}|,|\nabla^j\operatorname{Rm}|.
\]
For \(m=1\), apply this and the corresponding inequality for \(|\operatorname{Rm}|^2\)
to \(t|\nabla\operatorname{Rm}|^2+\beta|\operatorname{Rm}|^2\), choosing
\(\beta=\beta(n,\max\{\alpha,1\})\) so that all mixed derivative terms are
absorbed.  The result is
\[
(\partial_t-\Delta)F\le C(n)\beta K^3,
\]
and the maximum principle gives \(F\le C(n,\alpha)K^2\).
Inductively, add suitable positive multiples of
\(t^j|\nabla^j\operatorname{Rm}|^2\), \(0\le j<m\), to
\(t^m|\nabla^m\operatorname{Rm}|^2\).  Young's inequality absorbs the unique top-order
term into \(-2t^m|\nabla^{m+1}\operatorname{Rm}|^2\), while the induction hypothesis bounds
all remaining products by \(C K^2\).  The maximum principle then yields
\(t^m|\nabla^m\operatorname{Rm}|^2\le C K^2\), which is the assertion.
\end{proof}

\begin{lemma}[Doubling-time estimate]
\label{lem:chow-ii-curvature-doubling-time}
\source{chow-et-al-ricci-flow-part-ii:page-0260}
\dcref{ch14:14.6}
\group{Curvature Derivative Estimates}

If a complete bounded-curvature Ricci flow satisfies \(|\operatorname{Rm}(g(0))|\le K\),
then \(|\operatorname{Rm}(g(t))|\le2K\) for \(0\le t\le(16K)^{-1}\).
\end{lemma}
\begin{proof}
\sketch The curvature evolution and Kato's inequality give, in the barrier sense,
\((\partial_t-\Delta)|\operatorname{Rm}|\le8|\operatorname{Rm}|^2\).  The maximum principle therefore
compares \(u(t)=\sup_M|\operatorname{Rm}(g(t))|\) with the solution of \(v'=8v^2\),
\(v(0)=K\).  Thus \(u(t)\le K/(1-8Kt)\), which is at most \(2K\) on the
stated interval.
\end{proof}

\begin{corollary}[Derivative estimates from an initial curvature bound]
\label{cor:chow-ii-derivative-estimates-initial-curvature-bound}
\source{chow-et-al-ricci-flow-part-ii:page-0260}
\uses{thm:chow-ii-global-curvature-derivative-estimates, lem:chow-ii-curvature-doubling-time}
\dcref{ch14:14.7}
\group{Curvature Derivative Estimates}


For \(m,n\in\mathbb N\), there is \(C(m,n)\) such that, whenever the maximum
principle applies and \(|\operatorname{Rm}(g(0))|\le K\),
\[
|\nabla^m\operatorname{Rm}(x,t)|\le\frac{C(m,n)K}{t^{m/2}}
\qquad\left(0<t\le\frac1{16K}\right).
\]
\end{corollary}
\begin{proof}
\sketch Lemma~\ref{lem:chow-ii-curvature-doubling-time} bounds the curvature by
\(2K\) on the stated interval.  Apply
Theorem~\ref{thm:chow-ii-global-curvature-derivative-estimates} with this
bound and a fixed dimensionless time parameter.
\end{proof}

\begin{remark}[Derivatives and scaling]
\label{rem:chow-ii-curvature-derivative-scaling}
\source{chow-et-al-ricci-flow-part-ii:page-0260,chow-et-al-ricci-flow-part-ii:page-0261}
\dcref{ch14:14.8}
\group{Scaling}

Under \(g\mapsto cg\), distances scale by \(c^{1/2}\), while
\[
|\operatorname{Rm}|\sim c^{-1},\qquad |\nabla^m\operatorname{Rm}|\sim c^{-(m+2)/2}.
\]
Under parabolic rescaling \(\widetilde g(\widetilde t)=cg(t)\),
\(\widetilde t=ct\), the quantity
\(t^{m/2}|\nabla^m\operatorname{Rm}|\) scales in the same way as \(|\operatorname{Rm}|\).  Hence the
estimate in Theorem~\ref{thm:chow-ii-global-curvature-derivative-estimates}
is scale-covariant.
\end{remark}

\begin{remark}[Smoothing character of the derivative estimates]
\label{rem:chow-ii-derivative-estimates-smoothing}
\source{chow-et-al-ricci-flow-part-ii:page-0261}
\uses{thm:chow-ii-global-curvature-derivative-estimates}
\dcref{ch14:14.9}
\group{Curvature Derivative Estimates}


The factor \(t^{-m/2}\) permits no derivative bound at time zero, but gives
bounds of every order at each positive time.  For the heat equation this rate
is sharp: the solutions \(e^{-k^2t}\sin(kx)\) on the circle have uniformly
bounded size and gradients of order \(t^{-1/2}\) when \(k^2t\) is fixed.
\end{remark}

\section{Shi's local derivative estimates}

\begin{theorem}[Local first derivative estimate, Shi]
\label{thm:chow-ii-shi-local-first-derivative-estimate}
\source{chow-et-al-ricci-flow-part-ii:page-0263,chow-et-al-ricci-flow-part-ii:page-0264,chow-et-al-ricci-flow-part-ii:page-0265,chow-et-al-ricci-flow-part-ii:page-0266}
\uses{lem:chow-ii-cutoff-gradient-laplacian-estimates}
\dcref{ch14:14.10}
\group{Local Curvature Estimates}


Let \(\alpha>0\).  Suppose \(\overline B_{g(0)}(p,r)\Subset\mathcal U\),
\(0<\tau\le\alpha/K\), and \(|\operatorname{Rm}|\le K\) on
\(\mathcal U\times[0,\tau]\).  Then
\[
|\nabla\operatorname{Rm}(y,t)|\le
\frac{C(\alpha,K,r,n)}{\sqrt t}
=\frac{C(\alpha,\sqrt K r,n)K}{\sqrt t}
\]
on \(B_{g(0)}(p,r/2)\times(0,\tau]\).  If additionally
\(\gamma K^{-1/2}\le r\le\beta K^{-1/2}\), the last constant may be chosen
to depend only on \(n,\alpha,\beta,\gamma\).
\end{theorem}
\begin{proof}
\sketch The standard curvature evolution inequalities imply, for
\[
F=(16K^2+|\operatorname{Rm}|^2)|\nabla\operatorname{Rm}|^2,
\]
that
\[
(\partial_t-\Delta)F\le-\frac{F^2}{289K^4}+C(n)K^6.
\]
Indeed, the mixed term is bounded using
\(-\tfrac12a^4+8Ka^2b-32K^2b^2\le0\), and the lower-order curvature term is
absorbed by another half of \(-a^4\).  For \(G=tF\) and the cutoff \(\eta\)
from Lemma~\ref{lem:chow-ii-cutoff-gradient-laplacian-estimates}, set
\(H=\eta G\).  At a positive maximum of \(H\), the cutoff estimates and
\(t\le\alpha/K\) give
\[
0\le-\frac c{K^4}H^2
+\frac{C(\alpha,n)}{K^{5/2}r}H^{3/2}
+\frac{C(\alpha,\sqrt K r,n)}{r^2K}H+C(\alpha,n)K^4.
\]
The negative quadratic term dominates, so
\(H\le C(\alpha,\sqrt K r,n)K^4\).  Since \(\eta=1\) on the half ball and
\(16K^2t|\nabla\operatorname{Rm}|^2\le G\), this proves the first estimate.  Under the
two-sided restriction on \(r\sqrt K\), every coefficient in the displayed
polynomial is bounded in terms of \(n,\alpha,\beta,\gamma\), giving the final
claim.
\end{proof}

\begin{remark}[Independence from initial derivative bounds]
\label{rem:chow-ii-local-first-estimate-initial-data}
\source{chow-et-al-ricci-flow-part-ii:page-0263,chow-et-al-ricci-flow-part-ii:page-0264}
\uses{thm:chow-ii-shi-local-first-derivative-estimate}
\dcref{ch14:14.11}
\group{Local Curvature Estimates}


The estimate requires no bound on \(|\nabla\operatorname{Rm}(g(0))|\).  It is enough that
each smaller closed initial ball be compact; the flow need not be complete
outside the neighborhood on which the curvature bound is assumed.
\end{remark}

\begin{remark}[Dimensionless-radius form of Shi's estimate]
\label{rem:chow-ii-shi-estimate-dimensionless-radius}
\source{chow-et-al-ricci-flow-part-ii:page-0266}
\uses{thm:chow-ii-shi-local-first-derivative-estimate}
\dcref{ch14:14.12}
\group{Local Curvature Estimates}


If \(\gamma\le r\sqrt K\le\beta\), the maximum-principle polynomial in the
proof is dimensionless after dividing \(H\) by \(K^4\); consequently
\(|\nabla\operatorname{Rm}|\le C(\alpha,\beta,\gamma,n)K/\sqrt t\).
\end{remark}

\begin{theorem}[Interior first derivative estimate]
\label{thm:chow-ii-interior-first-derivative-estimate}
\source{chow-et-al-ricci-flow-part-ii:page-0267}
\dcref{ch14:14.13}
\group{Local Curvature Estimates}

If \(\overline B_{g(0)}(p,r)\Subset\mathcal U\) and \(|\operatorname{Rm}|\le K\) on
\(B_{g(0)}(p,r)\times[0,\tau]\), then
\[
|\nabla\operatorname{Rm}(y,t)|\le C(n)K
\left(r^{-2}+t^{-1}+K\right)^{1/2}
\]
for \(y\in B_{g(0)}(p,r/2)\) and \(0<t\le\tau\).
\end{theorem}
\begin{proof}
\sketch Fix \((y,t)\) and put \(A=r^{-2}+t^{-1}+K\).  After the parabolic rescaling
by \(A\), the curvature is at most one, the available spatial radius is at
least one, and the backward time interval has length at least a dimensional
constant.  The localized Bernstein quantity
\((16+|\operatorname{Rm}|^2)|\nabla\operatorname{Rm}|^2\), multiplied by a distance barrier vanishing on
the parabolic boundary, satisfies the same negative-quadratic inequality as
in Theorem~\ref{thm:chow-ii-shi-local-first-derivative-estimate}.  The maximum
principle bounds the rescaled gradient by \(C(n)\).  Scaling back multiplies
the bound by \(K\sqrt A\), proving the result.
\end{proof}

\begin{theorem}[Local higher derivative estimates, Shi]
\label{thm:chow-ii-shi-local-higher-derivative-estimates}
\source{chow-et-al-ricci-flow-part-ii:page-0267,chow-et-al-ricci-flow-part-ii:page-0268,chow-et-al-ricci-flow-part-ii:page-0269}
\uses{lem:chow-ii-shi-higher-derivative-quantity}
\dcref{ch14:14.14}
\group{Local Curvature Estimates}


For \(\alpha,K,r>0\) and \(m,n\in\mathbb N\), there is
\(C(\alpha,K,r,m,n)\) such that, under the hypotheses of
Theorem~\ref{thm:chow-ii-shi-local-first-derivative-estimate},
\[
|\nabla^m\operatorname{Rm}(y,t)|\le\frac{C(\alpha,K,r,m,n)}{t^{m/2}}
\]
on \(B_{g(0)}(p,r/2)\times(0,\tau]\).
\end{theorem}
\begin{proof}
\sketch The case \(m=1\) is
Theorem~\ref{thm:chow-ii-shi-local-first-derivative-estimate}.  Assume the
claim through order \(m\), first on successively smaller balls.  Apply
Lemma~\ref{lem:chow-ii-shi-higher-derivative-quantity} and multiply its
quantity \(F_m\) by a cutoff from
Lemma~\ref{lem:chow-ii-cutoff-bounded-derivatives}.  At a positive maximum of
\(G_m=\eta F_m\),
\[
0\le-cG_m^2+C+C(\alpha,K,r,m,n)G_m,
\]
so \(G_m\) is bounded.  Since \(\eta=1\) on the smaller ball, the definition
of \(F_m\) gives
\(t^{m+1}|\nabla^{m+1}\operatorname{Rm}|^2\le C\).  Induction proves the estimate on
\(B_{g(0)}(p,r/2^{m+1})\).  Choosing nested cutoff radii between \(r\) and
\(r/2\), instead of repeatedly halving the radius, gives the asserted common
half ball.
\end{proof}

\begin{lemma}[Evolution inequality for Shi's higher derivative quantity]
\label{lem:chow-ii-shi-higher-derivative-quantity}
\source{chow-et-al-ricci-flow-part-ii:page-0268,chow-et-al-ricci-flow-part-ii:page-0269,chow-et-al-ricci-flow-part-ii:page-0270}
\dcref{ch14:14.15}
\group{Local Curvature Estimates}

Suppose for \(1\le j\le m\) that
\[
|\nabla^j\operatorname{Rm}|\le C_jt^{-j/2}
\]
on \(B_{g(0)}(p,r/2^m)\times(0,\tau]\).  For a sufficiently large constant
\(C\), the quantity
\[
F_m=(C+t^m|\nabla^m\operatorname{Rm}|^2)t^{m+1}|\nabla^{m+1}\operatorname{Rm}|^2
\]
satisfies
\[
(\partial_t-\Delta)F_m\le-\frac c tF_m^2+\frac C t.
\]
\end{lemma}
\begin{proof}
\sketch For every \(k\), commuting covariant derivatives in the curvature equation
gives
\[
(\partial_t-\Delta)|\nabla^k\operatorname{Rm}|^2
\le-2|\nabla^{k+1}\operatorname{Rm}|^2
+C\sum_{i=0}^k|\nabla^i\operatorname{Rm}|,|\nabla^{k-i}\operatorname{Rm}|,|\nabla^k\operatorname{Rm}|.
\]
After multiplication by \(t^k\), the induction bounds control every product
which does not contain \(\nabla^{m+1}\operatorname{Rm}\) or \(\nabla^{m+2}\operatorname{Rm}\).  Applying
the product rule to \(F_m\), choosing
\(C\ge4t^m|\nabla^m\operatorname{Rm}|^2\), and estimating its gradient cross term reduces
the top-order terms to
\[
-10x^2+8xy-2y^2\le-\frac25y^2.
\]
Thus
\[
(\partial_t-\Delta)F_m
\le-\frac{2}{5t}\bigl(t^{m+1}|\nabla^{m+1}\operatorname{Rm}|^2\bigr)^2
+C(1+t)t^m|\nabla^{m+1}\operatorname{Rm}|^2+Ct+C/t.
\]
Increasing \(C\) and applying Young's inequality to the middle term yields
the stated inequality.
\end{proof}

\section{Modified local estimates and compactness}

\begin{theorem}[Shi's estimates with initial derivative bounds]
\label{thm:chow-ii-shi-estimates-initial-derivative-bounds}
\source{chow-et-al-ricci-flow-part-ii:page-0271,chow-et-al-ricci-flow-part-ii:page-0272,chow-et-al-ricci-flow-part-ii:page-0273}
\uses{lem:chow-ii-modified-shi-derivative-quantity}
\dcref{ch14:14.16}
\group{Local Curvature Estimates}


Let \(\ell\ge0\), \(n\ge2\), \(m\ge1\), and let
\(\alpha,K,K_\ell,r>0\).  Suppose \(0<\tau\le\alpha/K\),
\(\overline B_{g(0)}(p,r)\Subset\mathcal U\), and
\[
|\operatorname{Rm}(x,t)|\le K,
\qquad |\nabla^\beta\operatorname{Rm}(x,0)|\le K_\ell\quad(0\le\beta\le\ell)
\]
on the indicated ball and time interval.  Then
\[
|\nabla^m\operatorname{Rm}(y,t)|\le
\frac{C(\alpha,K,K_\ell,r,\ell,m,n)}{t^{(m-\ell)_+/2}}
\]
on \(B_{g(0)}(p,r/2)\times(0,\tau]\).  In particular, derivatives of order
at most \(\ell\) are uniformly bounded through time zero.
\end{theorem}
\begin{proof}
\sketch Induct on \(m\), using the initial bounds when the maximum of the localized
quantity occurs at time zero.  Lemma~\ref{lem:chow-ii-modified-shi-derivative-quantity}
and a cutoff from Lemma~\ref{lem:chow-ii-cutoff-bounded-derivatives} show at a
positive maximum of \(\eta F_{m,\ell}\) that either
\[
0\le-c(\eta F_{m,\ell})^2+C\eta^2+C\eta F_{m,\ell}
\]
or the same inequality after multiplication by \(t\).  Hence
\(\eta F_{m,\ell}\le C\), which gives
\[
|\nabla^{m+1}\operatorname{Rm}|^2\le Ct^{-(m-\ell+1)_+}
\]
on the next inner ball.  Nested radii approaching \(r/2\) complete the
induction and give the stated estimate.
\end{proof}

\begin{remark}[Recovery of Shi's estimates]
\label{rem:chow-ii-modified-shi-zero-order-case}
\source{chow-et-al-ricci-flow-part-ii:page-0271}
\uses{thm:chow-ii-shi-estimates-initial-derivative-bounds, thm:chow-ii-shi-local-higher-derivative-estimates}
\dcref{ch14:14.17}
\group{Local Curvature Estimates}


For \(\ell=0\), Theorem~\ref{thm:chow-ii-shi-estimates-initial-derivative-bounds}
has the same time powers as Shi's local derivative estimates.
\end{remark}

\begin{lemma}[Evolution inequality for the modified derivative quantity]
\label{lem:chow-ii-modified-shi-derivative-quantity}
\source{chow-et-al-ricci-flow-part-ii:page-0272,chow-et-al-ricci-flow-part-ii:page-0273,chow-et-al-ricci-flow-part-ii:page-0274,chow-et-al-ricci-flow-part-ii:page-0275,chow-et-al-ricci-flow-part-ii:page-0276,chow-et-al-ricci-flow-part-ii:page-0277}
\dcref{ch14:14.18}
\group{Local Curvature Estimates}

Fix \(\ell\ge0\) and suppose, for \(1\le j\le m\), that
\[
|\nabla^j\operatorname{Rm}|\le Ct^{-(j-\ell)_+/2}.
\]
Put
\[
F_{m,\ell}=
(\widetilde C+t^{(m-\ell)_+}|\nabla^m\operatorname{Rm}|^2)
t^{(m-\ell+1)_+}|\nabla^{m+1}\operatorname{Rm}|^2.
\]
For \(\widetilde C\) sufficiently large,
\[
(\partial_t-\Delta)F_{m,\ell}
\le-\frac c{t^{\operatorname{sgn}((m-\ell+1)_+)}}F_{m,\ell}^2
+\frac C{t^{\operatorname{sgn}((m-\ell+1)_+)}}.
\]
\end{lemma}
\begin{proof}
\sketch Set \(a=(m-\ell)_+\) and \(b=(m-\ell+1)_+\).  Apply the commuted curvature
inequality used in Lemma~\ref{lem:chow-ii-shi-higher-derivative-quantity} to
\(t^a|\nabla^m\operatorname{Rm}|^2\) and \(t^b|\nabla^{m+1}\operatorname{Rm}|^2\).  The induction
hypothesis controls every lower-order product because
\[
(i-\ell)_++(m+1-i-\ell)_+\le(m+1-\ell)_+.
\]
Choose \(\widetilde C\ge4t^a|\nabla^m\operatorname{Rm}|^2\).  In the product evolution,
the gradient cross term is absorbed by
\(-10x^2+8xy-2y^2\le-2y^2/5\).  Consequently
\[
(\partial_t-\Delta)F_{m,\ell}
\le-c,t^{a+b}|\nabla^{m+1}\operatorname{Rm}|^4
+Ct^a|\nabla^{m+1}\operatorname{Rm}|^2+C.
\]
If \(b=0\), Young's inequality gives \(-cF_{m,\ell}^2+C\).  If \(b>0\),
then \(b-a=1\), so factoring \(t^{-1}\) gives
\(-ct^{-1}F_{m,\ell}^2+Ct^{-1}\).  These are exactly the two cases in the
statement.
\end{proof}

\begin{theorem}[Modified Cheeger--Gromov compactness]
\label{thm:chow-ii-modified-cheeger-gromov-compactness}
\source{chow-et-al-ricci-flow-part-ii:page-0277}
\uses{thm:chow-ii-shi-estimates-initial-derivative-bounds}
\dcref{ch14:14.19}
\group{Ricci Flow Compactness}


Let \((M_k,g_k(t),p_k)\), \(t\in[0,T)\), be pointed complete Ricci flows.
Suppose that, for some \(m_0\in\mathbb N\cup\{\infty\}\), the initial
curvature derivatives through order \(m_0+1\) are uniformly bounded, the
space-time curvature is uniformly bounded on every \([0,T_0]\) with
\(T_0<T\), and \(\operatorname{inj}_{g_k(0)}(p_k)\ge\delta>0\).  Then a
subsequence converges in pointed \(C^{m_0}\) Cheeger--Gromov topology to a
pointed complete Ricci flow on \([0,T)\).
\end{theorem}
\begin{proof}
\sketch Initial derivative and injectivity-radius bounds give pointed
\(C^{m_0}\) Cheeger--Gromov convergence of a subsequence at time zero.
Theorem~\ref{thm:chow-ii-shi-estimates-initial-derivative-bounds}, applied on
each compact time interval, supplies uniform space-time bounds through order
\(m_0+1\).  The Ricci flow equation converts these to the corresponding time
derivative bounds.  Arzela--Ascoli and a diagonal exhaustion therefore give
\(C^{m_0}\) convergence on compact subsets of space-time.  Passing to the
limit in \(\partial_tg=-2\operatorname{Rc}\) yields a Ricci flow.  Uniform equivalence on
compact time intervals and completeness of every \(g_k(t)\) imply completeness
of the limit.
\end{proof}

\begin{remark}[Finite-regularity compactness]
\label{rem:chow-ii-finite-regularity-compactness}
\source{chow-et-al-ricci-flow-part-ii:page-0277}
\uses{thm:chow-ii-modified-cheeger-gromov-compactness}
\dcref{ch14:14.20}
\group{Ricci Flow Compactness}


The theorem retains a prescribed finite degree of regularity at the initial
time; no positive-time smoothing estimate alone supplies this conclusion at
\(t=0\).
\end{remark}

\section{Applications}

\begin{lemma}[Existence of curvature bumps]
\label{lem:chow-ii-existence-curvature-bumps}
\source{chow-et-al-ricci-flow-part-ii:page-0278,chow-et-al-ricci-flow-part-ii:page-0279}
\uses{thm:chow-ii-interior-first-derivative-estimate}
\dcref{ch14:14.21}
\group{Curvature Positivity}


Suppose \(|\operatorname{Rm}|\le K\) on \(B_{g(0)}(p,r)\times[0,T]\) and the least
curvature-operator eigenvalue satisfies
\(\lambda_1(\operatorname{Rm})(p,\tau)\ge cK\).
\begin{enumerate}
\item There is \(\eta=\eta(n,K,r,\tau,c)>0\) such that
\(\lambda_1(\operatorname{Rm})(x,\tau)\ge cK/2\) on \(B_{g(0)}(p,\eta)\).
\item If \(r\ge cK^{-1/2}\) and \(\tau\ge cK^{-1}\), there is
\(\gamma(n,c)>0\) such that the same conclusion holds on
\(B_{g(0)}(p,\gamma K^{-1/2})\).
\end{enumerate}
\end{lemma}
\begin{proof}
\sketch The local derivative estimate bounds \(|\nabla\operatorname{Rm}(\cdot,\tau)|\) on the half
ball.  Along a minimizing \(g(\tau)\)-geodesic \(\beta\) from \(p\) to \(x\),
\[
\lambda_1(\operatorname{Rm})(x,\tau)
\ge cK-\int_\beta|\nabla\operatorname{Rm}(\cdot,\tau)|\,ds_{g(\tau)}.
\]
Metric comparison and a sufficiently small \(g(0)\)-radius make the integral
at most \(cK/2\), proving the first assertion.  Under the hypotheses of the
second assertion,
Theorem~\ref{thm:chow-ii-interior-first-derivative-estimate} gives
\(|\nabla\operatorname{Rm}|\le C(n,c)K^{3/2}\).  Taking
\(\gamma\le\min\{c/(2C(n,c)),c/2\}\) proves the uniform-radius claim.
\end{proof}

\begin{lemma}[Instantaneous local derivative bounds]
\label{lem:chow-ii-instantaneous-local-derivative-bounds}
\source{chow-et-al-ricci-flow-part-ii:page-0279}
\uses{thm:chow-ii-interior-first-derivative-estimate}
\dcref{ch14:14.22}
\group{Local Curvature Estimates}


Let \(g(t)\), \(t\in(0,T)\), be complete with bounded nonnegative curvature
operator.  If
\(R(x,\tau)\le K\) on \(B_{g(\tau)}(p,r)\), then
\[
|\nabla\operatorname{Rm}(x,\tau)|\le C(n)K(r^{-2}+\tau^{-1}+K)^{1/2}
\]
on \(B_{g(\tau)}(p,r/2)\).
\end{lemma}
\begin{proof}
\sketch Hamilton's trace Harnack inequality makes \(tR(x,t)\) nondecreasing, hence
\(R(x,t)\le2K\) for \(t\in[\tau/2,\tau]\) on the terminal ball.  Nonnegative
curvature operator converts this to \(|\operatorname{Rm}|\le C(n)K\).  Nonnegative Ricci
curvature makes distances nonincreasing, so the radius-\(r\) ball at time
\(\tau/2\) lies in the terminal ball.  Apply
Theorem~\ref{thm:chow-ii-interior-first-derivative-estimate} on
\([\tau/2,\tau]\).
\end{proof}

\section{Local heat equation and local Ricci flow}

\begin{lemma}[Bernstein estimate for the local heat equation]
\label{lem:chow-ii-local-heat-bernstein-estimate}
\source{chow-et-al-ricci-flow-part-ii:page-0280,chow-et-al-ricci-flow-part-ii:page-0281}
\dcref{ch14:14.23}
\group{Local Heat Equation}

Let \(\chi\) be smooth and compactly supported in a bounded domain \(\Omega\)
of a complete Riemannian manifold with \(\operatorname{Rc}\ge-Kg\).  If
\(\partial_tu=\chi^2\Delta u\) and \(u(0)=u_0\), then
\[
|\nabla u|^2\le\frac C{\chi^2t}
\]
on \(\Omega\times(0,T]\), where \(C\) depends only on \(n,K,T\), the size
of \(u_0\), and upper bounds for \(-\chi\Delta\chi\), \(|\nabla\chi|^2\),
and \(\chi^2\).
\end{lemma}
\begin{proof}
\sketch Direct calculation gives
\[
(\partial_t-\chi^2\Delta)u^2=-2\chi^2|\nabla u|^2
\]
and, using \(|\nabla^2u|^2\ge n^{-1}(\Delta u)^2\),
\[
(\partial_t-\chi^2\Delta)(\chi^2|\nabla u|^2)
\le2(-\chi\Delta\chi+(7+2n)|\nabla\chi|^2+K\chi^2)
\chi^2|\nabla u|^2.
\]
If all three coefficient functions are bounded above by \(C_0\), then
\[
(\partial_t-\chi^2\Delta)
\left(\chi^2t|\nabla u|^2+\bigl((8+2n+K)C_0T+\tfrac12\bigr)u^2\right)
\le0.
\]
The first summand vanishes on the initial and lateral parabolic boundary.
The maximum principle bounds the displayed quantity by its initial maximum,
and division by \(\chi^2t\) proves the estimate.
\end{proof}

\begin{example}[Higher derivative estimates for the local heat equation]
\label{ex:chow-ii-local-heat-higher-derivative-estimates}
\source{chow-et-al-ricci-flow-part-ii:page-0281}
\uses{lem:chow-ii-local-heat-bernstein-estimate}
\dcref{ch14:14.24}
\group{Local Heat Equation}


For every \(m\ge1\), repeated differentiation of
\(\partial_tu=\chi^2\Delta u\), followed by the Bernstein argument of
Lemma~\ref{lem:chow-ii-local-heat-bernstein-estimate}, yields on
\(\{\chi>0\}\times(0,T]\) an estimate
\[
|\nabla^mu|^2\le C_m\chi^{-2m}t^{-m},
\]
where \(C_m\) depends on curvature bounds and on derivatives of \(\chi\) and
\(u_0\) through the orders occurring after \(m\) differentiations.
\end{example}

\begin{proposition}[Short-time existence for local Ricci flow]
\label{prop:chow-ii-local-ricci-flow-short-time-existence}
\source{chow-et-al-ricci-flow-part-ii:page-0282,chow-et-al-ricci-flow-part-ii:page-0283}
\dcref{ch14:14.25}
\group{Local Ricci Flow}

Let \((M,g_0)\) be complete and let \(\chi\) be smooth with compact support.
There is \(\delta>0\), depending only on \(g_0\) near
\(\operatorname{supp}\chi\), and a solution of
\[
\partial_tg=-2\chi^2\operatorname{Rc}(g),\qquad g(0)=g_0,
\]
on \([0,\delta)\).
\end{proposition}
\begin{proof}
\sketch Fix a background connection and solve the local Ricci--DeTurck equation
\[
\partial_tg_{ij}=-2\chi^2R_{ij}
+\nabla_i(\chi^2W)_j+\nabla_j(\chi^2W)_i,
\qquad
W_j=g_{jk}g^{pq}(\Gamma^k_{pq}-\widetilde\Gamma^k_{pq}).
\]
Its linearization has principal part \(\chi^2\Delta h\); outside the compact
support it is the stationary equation, and on the support standard local
quasilinear parabolic existence gives a solution for a time depending only on
the initial coefficients there.  The compactly supported vector field
\(-\chi^2W\) integrates to diffeomorphisms \(\varphi_t\).  Differentiating a
pullback shows
\(\partial_t(\varphi_t^*g)=-2\chi^2\operatorname{Rc}(\varphi_t^*g)\), and
\(\varphi_0^*g_0=g_0\), which proves the assertion.
\end{proof}

\begin{theorem}[Existence-time estimate for local Ricci flow]
\label{thm:chow-ii-local-ricci-flow-existence-time}
\source{chow-et-al-ricci-flow-part-ii:page-0283}
\uses{prop:chow-ii-local-ricci-flow-short-time-existence}
\dcref{ch14:14.26}
\group{Local Ricci Flow}


There is \(C(n)>0\) such that, if \((M^n,g_0)\) is complete and
\(|\operatorname{Rm}(g_0)|\le K\), the local Ricci flow exists at least for
\[
0\le t\le C(n)\bigl(\sup_M|\nabla\chi|^2+K\bigr)^{-1}.
\]
\end{theorem}
\begin{proof}
\sketch Parabolically rescale so that \(\sup|\nabla\chi|^2+K=1\).  The local
Ricci--DeTurck system in the preceding proof then has uniformly bounded
initial principal coefficients and lower-order coefficients on the compact
support of \(\chi\).  Its standard interior energy estimates give a
dimension-only existence interval \([0,C(n)]\); the continuation criterion
shows that the solution cannot cease earlier while those estimates hold.
Undoing the parabolic rescaling gives the claimed time.  Pulling back by the
DeTurck diffeomorphisms preserves the interval.
\end{proof}

\begin{lemma}[Variation formula for the Riemann tensor]
\label{lem:chow-ii-riemann-tensor-variation-formula}
\source{chow-et-al-ricci-flow-part-ii:page-0283,chow-et-al-ricci-flow-part-ii:page-0284,chow-et-al-ricci-flow-part-ii:page-0285}
\dcref{ch14:14.27}
\group{Curvature Variations}

If \(\partial_sg_{ij}=-2h_{ij}\), then
\[
\partial_sR_{ijkl}
=(\nabla^2\mathbin\odot h)_{ijkl}
-\frac12\left(
R_{qjkl}h_{iq}+R_{iqkl}h_{jq}+R_{ijq\ell}h_{kq}+R_{ijkq}h_{\ell q}
\right),
\]
where \(\nabla^2_{ij}=\tfrac12(\nabla_i\nabla_j+\nabla_j\nabla_i)\) and
\(\odot\) is the Kulkarni--Nomizu product.
\end{lemma}
\begin{proof}
\sketch Differentiating the Levi--Civita connection and then the curvature gives
\[
\partial_sR_{ijkl}
=\nabla_i\nabla_\ell h_{jk}+\nabla_j\nabla_kh_{i\ell}
-\nabla_i\nabla_kh_{j\ell}-\nabla_j\nabla_\ell h_{ik}
-R^q{}_{ijk}h_{q\ell}+R^q{}_{ij\ell}h_{kq}.
\]
Commute each pair of derivatives to replace it by its symmetrization.  The
commutators contribute curvature contracted with \(h\).  The first Bianchi
identity combines these contractions into the four terms displayed in the
statement, while the four symmetrized Hessian terms are precisely
\(\nabla^2\mathbin\odot h\).
\end{proof}
